\documentclass[11pt,a4paper,oneside]{report}
\usepackage[utf8]{inputenc}
\usepackage[T1]{fontenc}
\usepackage[italian]{babel}
\usepackage[margin=2.5cm,headheight=14pt]{geometry}
\usepackage{graphicx}
\usepackage{array,tabularx,longtable,booktabs}
\usepackage{enumitem}
\usepackage{xcolor}
\usepackage{xurl}
\usepackage{hyperref}
\usepackage{fancyhdr}
\hypersetup{colorlinks=true,linkcolor=black,urlcolor=blue!50!black,
  pdftitle={BugBoard26 - Documentazione finale},pdfauthor={Gruppo BugBoard26}}
\pagestyle{fancy}
\fancyhf{}
\fancyhead[L]{BugBoard26}
\fancyhead[R]{Documentazione finale}
\fancyfoot[C]{\thepage}
\setlength{\parindent}{0pt}
\setlength{\parskip}{5pt}
\setlength{\emergencystretch}{3em}
\setlist{itemsep=2pt,topsep=4pt}
\renewcommand{\arraystretch}{1.18}
\newcolumntype{P}[1]{>{\raggedright\arraybackslash}p{#1}}
\newcommand{\diagramma}[2]{%
  \begin{figure}[htbp]\centering
  \includegraphics[width=\textwidth,height=0.76\textheight,keepaspectratio]{#1}
  \caption{#2}\end{figure}}

\begin{document}
\hypersetup{pageanchor=false}
\begin{titlepage}
\centering
\vspace*{2cm}
{\Huge\bfseries BugBoard26\par}
\vspace{0.6cm}
{\LARGE Documentazione finale\par}
\vspace{0.6cm}
{\large Requisiti, design e testing\par}
\vspace{1cm}
Progetto di Ingegneria del Software\\A.A. 2025/2026 -- Canale 2
\vfill
\begin{tabular}{ll}
Componenti & Valerio Porcile -- N86004517\\
 & Emanuel Giuseppe Lupoli -- N86004280\\
 & Emanuele Martucci -- N86004084\\
Modalita' & Tradizionale\\
Versione & 1.6 -- relazione tecnica\\
Data di riferimento & 19 ottobre 2026\\
Codice descritto (commit Git) & \texttt{b95d008}\\
\end{tabular}
\vfill
\end{titlepage}
\pagenumbering{roman}
\hypersetup{pageanchor=true}

\chapter*{Introduzione}
\addcontentsline{toc}{chapter}{Introduzione}
Questo documento riunisce la specifica dei requisiti, il design del sistema
e le attivita' di testing in capitoli distinti, secondo la sezione 8 della
traccia ufficiale v1.1. Riprende il materiale dei due homework e lo confronta
con il codice disponibile nel repository.

La relazione presenta il perimetro assegnato, i casi d'uso, le scelte
architetturali e Object-Oriented, il modello dei dati e le verifiche svolte.
Il capitolo sul testing distingue test automatici, collaudo del sistema,
analisi di qualita' e valutazione dell'usabilita'. Il capitolo conclusivo
espone i limiti del prodotto e della valutazione.

La revisione del codice descritta e' \texttt{b95d008}. I risultati delle
verifiche mantengono il proprio riferimento storico: il collaudo del
16 ottobre riguarda \texttt{e139951}; l'audit del 19 ottobre usa la base
\texttt{6aa175a} con correzioni locali, mentre la successiva review frontend
usa \texttt{06ae1b0} con modifiche locali. I report conservano stato Git e
impronte dei sorgenti per identificare gli esemplari analizzati. La revisione
editoriale non costituisce una nuova esecuzione dei test sul commit di
copertina.

\textbf{Repository:} \url{https://github.com/Valeriop03/BugBoard26}.
La traccia di riferimento e' il file
\path{Progetto-SWE-2025-2026 - canale 2 - v1.1.pdf} nella radice.
I sorgenti dei homework sono conservati come materiale precedente.
I percorsi relativi citati nella relazione si riferiscono alla cartella
\path{Documentazione/Documentazione finale/}, salvo diversa indicazione.

{\small\setlength{\parskip}{0pt}\tableofcontents}
\clearpage
\pagenumbering{arabic}

\chapter{Specifica dei requisiti}
\section{Scopo e perimetro}
BugBoard26 permette a un gruppo di lavoro di segnalare e consultare issue,
seguirne lo stato e comunicare la risoluzione al segnalatore. Il progetto
e' una web app distribuita con frontend e backend indipendenti.
Le funzionalita' assegnate sono esclusivamente quelle riportate sotto.

\begin{longtable}{P{0.09\textwidth} P{0.84\textwidth}}
\toprule N. traccia & Funzionalita' assegnata\\\midrule\endhead
1 & Login email/password, amministratore iniziale e creazione delle utenze da parte di ADMIN.\\
2 & Creazione issue con titolo, descrizione, tipo, priorita' e immagine opzionali; stato iniziale TODO.\\
3 & Lista riepilogativa, filtri e ordinamento.\\
6 & Cambio stato da parte dell'assegnatario; notifica al segnalatore alla risoluzione.\\
8 & Esportazione dell'elenco in formato condivisibile, realizzata con CSV.\\
11 & Ricerca per parole chiave.\\
13 & Archiviazione riservata ad ADMIN e consultazione in una pagina dedicata.\\
14 & Suggerimento dell'assegnatario in base al carico di lavoro corrente.\\
15 & Utenze READONLY create dagli amministratori per la consultazione.\\
16 & Contrassegno duplicato riservato ad ADMIN e chiusura logica immediata.\\
\bottomrule
\end{longtable}

Non sono incluse gestione dei commenti, dashboard, tag, cronologia,
report mensili o scadenze. Il campo assegnatario e' necessario per il
cambio stato e il suggerimento, ma non introduce il flusso separato di
assegnazione e relativa notifica della funzionalita' 4, non assegnata.
La notifica di risoluzione e' interna all'applicazione, non un'email.

\clearpage
\section{Glossario e utenti target}
\begin{longtable}{P{0.23\textwidth} P{0.70\textwidth}}
\toprule Termine & Significato\\\midrule\endhead
Issue & Segnalazione di tipo QUESTION, BUG, DOCUMENTATION o FEATURE.\\
Segnalatore & Utente che ha creato la issue.\\
Assegnatario & Utente indicato nel campo \texttt{AssignedToId}.\\
Archiviata & Issue esclusa dalla lista principale, ma conservata e consultabile.\\
Duplicata & Issue chiusa logicamente come DUPLICATE e collegata a un originale.\\
Notifica & Messaggio persistente per un singolo destinatario, con stato di lettura.\\
\bottomrule
\end{longtable}

Gli utenti target sono membri di un team di sviluppo e referenti esterni.
I primi devono registrare e consultare segnalazioni durante il lavoro;
i secondi devono seguirne l'avanzamento senza modificare le issue.
Non e' previsto l'accesso anonimo ai dati del progetto.

\begin{tabularx}{\textwidth}{P{0.22\textwidth} X}
\toprule Attore & Responsabilita'\\\midrule
ADMIN & Crea utenti e issue, consulta i dati, cambia stato anche senza assegnazione, archivia e marca duplicati.\\
USER & Crea e consulta issue; cambia stato soltanto se assegnatario.\\
READONLY & Consulta issue, dettagli, immagini e archivio; usa ricerca, filtri ed export senza modificare issue.\\
\bottomrule
\end{tabularx}

Tutti i ruoli autenticati possono consultare il suggerimento assegnatario
e le proprie notifiche.

\textbf{Interpretazione adottata per RF-17:} READONLY puo'
consultare le proprie notifiche, ma non segnarle come lette. Questa scelta
applica il divieto di scrittura anche allo stato di lettura. Il contratto e' implementato e
verificato con test HTTP e browser: GET proprie 200, PATCH read 403 senza
scritture e nessun comando UI, anche per notifiche gia' lette.

\section{Requisiti funzionali}
Gli identificativi seguenti sono mantenuti dalla specifica dei requisiti;
non coincidono con i numeri delle funzionalita' nella traccia.
\begin{longtable}{P{0.12\textwidth} P{0.81\textwidth}}
\toprule ID & Requisito\\\midrule\endhead
RF-01 & Login tramite email e password.\\
RF-02 & Account amministratore iniziale gia' attivo.\\
RF-03 & Creazione utenti da ADMIN con email, password e ruolo.\\
RF-04 & Creazione issue da utente autenticato non READONLY con titolo e descrizione.\\
RF-05 & Tipo della issue, priorita' opzionale e immagine opzionale.\\
RF-06 & Stato iniziale TODO.\\
RF-07 & Lista delle issue non archiviate.\\
RF-08 & Filtri e ordinamento per criteri rilevanti.\\
RF-09 & Ricerca per parole chiave su titolo e descrizione.\\
RF-10 & Cambio stato da parte dell'assegnatario; ADMIN ha accesso completo.\\
RF-11 & Notifica al segnalatore quando la issue passa a RESOLVED.\\
RF-12 & Esportazione dell'elenco in CSV.\\
RF-13 & Archiviazione riservata ad ADMIN.\\
RF-14 & Esclusione degli archiviati dalla lista principale.\\
RF-15 & Consultazione delle issue archiviate in una pagina dedicata.\\
RF-16 & Suggerimento assegnatario basato sul carico corrente.\\
RF-17 & READONLY visualizza issue e dettagli senza modificare dati.\\
RF-18 & ADMIN puo' contrassegnare una issue come duplicata di un'altra.\\
RF-19 & Chiusura immediata della issue contrassegnata come duplicata.\\
\bottomrule
\end{longtable}

\section{Requisiti non funzionali e vincoli}
\begin{longtable}{P{0.12\textwidth} P{0.81\textwidth}}
\toprule ID & Requisito\\\midrule\endhead
RNF-01 & Password memorizzate tramite hashing, non in chiaro.\\
RNF-02 & Verifica dei permessi nelle API, non soltanto nell'interfaccia.\\
RNF-03 & Frontend e backend separati; comunicazione esclusivamente via API REST.\\
RNF-04 & Stato persistente gestito centralmente dal backend e dal database.\\
RNF-05 & Interfaccia leggibile, con operazioni disponibili coerenti con il ruolo.\\
RNF-06 & Lista consultabile fluidamente con i dati dimostrativi; nessuna prestazione su larga scala dichiarata senza misure.\\
RNF-07 & Separazione tra controller, servizi, modelli e accesso ai dati.\\
RNF-08 & Avvio locale ripetibile seguendo il README.\\
\bottomrule
\end{longtable}
Il codice usa linguaggi Object-Oriented e non delega logica e persistenza
a servizi MBaaS. La traccia richiede strumenti CASE per il design.
Il deployment su cloud e' auspicato, ma non obbligatorio: non costituisce
una funzionalita' mancante del progetto.

\section{Modello dei casi d'uso}
Il diagramma rappresenta i ruoli e le funzionalita' assegnate. Il login
precede le operazioni protette. La notifica e' legata alla risoluzione;
la chiusura del duplicato non richiede un comando separato.
\diagramma{../Primo homework/UCD/use_case_diagram.png}{Diagramma dei casi d'uso.}
\clearpage

\section{UC-01 -- Segnalare una issue}
\begin{tabularx}{\textwidth}{P{0.25\textwidth} X}
\toprule Campo & Descrizione\\\midrule
Attore e scopo & USER oppure ADMIN registra una nuova segnalazione.\\
Precondizioni & Account attivo e autenticato, non READONLY.\\
Postcondizione & Issue persistita in TODO e consultabile nella lista e nel dettaglio.\\
Scenario principale & 1. L'utente apre il form. 2. Inserisce titolo, descrizione e tipo; puo' scegliere priorita' e immagine. 3. Il sistema mostra il suggerimento assegnatario. 4. Valida i dati e crea la issue. 5. Se selezionata, carica l'immagine. 6. Apre il dettaglio della issue creata.\\
Alternative & Dati mancanti/non validi: nessuna creazione. Suggerimento indisponibile: creazione senza assegnatario. Upload fallito dopo creazione: issue conservata, possibilita' di riprovare solo l'upload o aprire il dettaglio senza immagine.\\
Requisiti & RF-04, RF-05, RF-06, RF-07, RF-16.\\
\bottomrule
\end{tabularx}
Il form richiede almeno tre caratteri dopo la rimozione degli spazi esterni
per titolo e descrizione;
i limiti massimi sono rispettivamente 200 e 4000 caratteri.
Il suggerimento viene usato come assegnatario dal form; non e' una nuova
API di assegnazione. Se indicato, l'assegnatario deve essere attivo e non
READONLY. La creazione JSON e l'upload multipart sono due richieste distinte.
\diagramma{../Primo homework/mock-up/Creazione issue.png}{Mock-up di creazione: riferimento progettuale, non screenshot del collaudo.}
\clearpage

\section{UC-02 -- Risolvere una issue assegnata}
\begin{tabularx}{\textwidth}{P{0.25\textwidth} X}
\toprule Campo & Descrizione\\\midrule
Attore e scopo & USER assegnatario oppure ADMIN aggiorna l'avanzamento della issue.\\
Precondizioni & Account attivo e autenticato; issue esistente e non contrassegnata come duplicata.\\
Postcondizioni & Nuovo stato salvato; alla risoluzione sono valorizzati data di risoluzione e notifica non letta per il segnalatore.\\
Scenario principale & 1. L'utente apre il dettaglio. 2. Seleziona RESOLVED. 3. Il backend verifica ruolo e assegnazione. 4. Aggiorna stato e date, crea la notifica e salva. 5. Il frontend mostra il risultato aggiornato.\\
Alternative & USER non assegnatario o READONLY: 403. Issue inesistente: 404. Stato gia' corrente: risposta senza nuova notifica o nuove date. Issue duplicata: 409. Token non valido: 401.\\
Requisiti & RF-10, RF-11, RF-17, RNF-02.\\
\bottomrule
\end{tabularx}
Non e' imposta una matrice aggiuntiva di transizioni. Il passaggio da
RESOLVED a un altro stato azzera la data di risoluzione; una risoluzione
successiva genera una nuova notifica. Le notifiche precedenti restano.
\diagramma{../Primo homework/mock-up/Dettaglio Issue.png}{Mock-up del dettaglio e del cambio stato.}
\clearpage

\section{UC-03 -- Contrassegnare un duplicato}
\begin{tabularx}{\textwidth}{P{0.25\textwidth} X}
\toprule Campo & Descrizione\\\midrule
Attore e scopo & ADMIN collega una segnalazione duplicata all'originale.\\
Precondizioni & Autenticazione ADMIN; entrambe le issue esistono e sono distinte.\\
Postcondizioni & Stato DUPLICATE, riferimento all'originale salvato, \texttt{ResolvedAt} nullo.\\
Scenario principale & 1. Apre il dettaglio del duplicato. 2. Inserisce l'ID dell'originale. 3. Il backend verifica i riferimenti. 4. Salva stato e collegamento. 5. Il frontend aggiorna il dettaglio.\\
Alternative & Originale assente o autoreferenza: 400. Issue corrente assente: 404. Ruolo non ADMIN: 403.\\
Requisiti & RF-18, RF-19, RNF-02.\\
\bottomrule
\end{tabularx}
DUPLICATE rappresenta la chiusura logica, non l'eliminazione. Il percorso
di cambio stato non permette di impostarlo direttamente: occorre l'API
dedicata con il riferimento all'originale.

\chapter{Design del sistema e del software}
\section{Architettura e scelte tecnologiche}
Il frontend Angular presenta i dati e invia richieste HTTP. ASP.NET Core
Web API espone le operazioni, controlla i permessi e gestisce la
persistenza tramite Entity Framework Core e PostgreSQL. Il backend
funziona anche senza frontend; nessuna pagina Angular accede direttamente
al database o alla directory delle immagini.

Angular e' usato per routing, form e servizi HTTP. ASP.NET Core permette
di organizzare endpoint, dependency injection e autorizzazione nello
stesso progetto. PostgreSQL viene avviato in Docker per ridurre le
differenze di configurazione fra i tre PC. Il progetto backend usa .NET 10,
il frontend Angular 19; le versioni effettive dei pacchetti sono nei
rispettivi file di progetto e lockfile.
\diagramma{../Secondo homework/diagrammi_architetturali/component_diagram.png}{Componenti e comunicazione del sistema.}
\diagramma{../Secondo homework/diagrammi_architetturali/deployment_locale.png}{Deployment locale con frontend, API e database.}
\clearpage

\section{Organizzazione Object-Oriented}
\begin{tabularx}{\textwidth}{P{0.35\textwidth} X}
\toprule Elemento & Responsabilita'\\\midrule
Controller & Gestiscono richieste e risposte HTTP, autorizzazioni e coordinamento delle operazioni.\\
\texttt{IssueDomainService} & Logiche semplici di ricerca in memoria, permessi sullo stato e suggerimento assegnatario.\\
\texttt{ApplicationDbContext} & Mapping, relazioni e accesso a PostgreSQL con Entity Framework Core.\\
\texttt{User}, \texttt{Issue}, \texttt{Notification} & Entita' del dominio con proprieta' e relazioni.\\
\texttt{IPasswordHasher}, \texttt{IJwtTokenService} & Contratti per hashing password e generazione dei token; implementazioni iniettate nei controller.\\
\texttt{IssueImageStorage} & Validazione e lettura/scrittura delle immagini in una directory privata.\\
DTO di richiesta/risposta & Separano i dati delle API dalle entita' persistenti; non espongono password hash o percorsi fisici.\\
\bottomrule
\end{tabularx}
La separazione rende testabili le logiche senza HTTP e consente di cambiare
la presentazione mantenendo le API. Le interfacce di autenticazione
separano il contratto dall'implementazione. Non sono introdotti repository
generici o livelli enterprise: per questo perimetro il DbContext e i
servizi presenti sono sufficienti. Alcune query e operazioni sono nei
controller; non si dichiara una separazione assoluta di tutta la logica.
\diagramma{../Secondo homework/diagrammi_dettaglio/class_diagram_parte1.png}{Classi di dominio, persistenza e servizi.}
\diagramma{../Secondo homework/diagrammi_dettaglio/class_diagram_parte2.png}{Controller e dipendenze.}
\clearpage

\section{Persistenza}
Gli identificativi sono interi. Gli enum sono memorizzati come stringhe;
le date applicative sono prodotte in UTC.
\begin{longtable}{P{0.17\textwidth} P{0.76\textwidth}}
\toprule Entita' & Campi principali\\\midrule\endhead
User & Id, Email, PasswordHash, Role, CreatedAt, IsActive. Email univoca; collezioni di issue create/assegnate e notifiche.\\
Issue & Id, Title, Description, Type, Priority nullable, Status, CreatedById, AssignedToId nullable, DuplicateOfIssueId nullable, ImagePath nullable, IsArchived, CreatedAt, UpdatedAt, ResolvedAt nullable.\\
Notification & Id, UserId, IssueId, Message, IsRead, CreatedAt.\\
\bottomrule
\end{longtable}
Ogni issue ha un segnalatore e puo' avere un assegnatario. La relazione
con l'originale del duplicato e' un autoriferimento opzionale. Le notifiche
appartengono a un destinatario e a una issue. Il mapping impedisce la
cancellazione del segnalatore o dell'originale referenziato; la cancellazione
dell'assegnatario azzera il riferimento. Non sono esposte nuove funzioni
di cancellazione per sfruttare queste regole.

I byte delle immagini non sono in PostgreSQL: il database conserva un nome
gestito dal backend. In Docker file e database sono mantenuti in due volumi
distinti. Il seed applica le migrazioni e crea l'ADMIN se non esiste gia'
un amministratore. Non reimposta la password di account esistenti.
\diagramma{../Secondo homework/diagrammi_dettaglio/database_schema.png}{Schema della persistenza.}
\clearpage

\section{Autenticazione e autorizzazione}
Il login normalizza l'email, verifica account attivo e password con
\texttt{Pbkdf2PasswordHasher}, quindi restituisce JWT, scadenza e identita'
dell'utente. Il frontend conserva la sessione in localStorage, ne controlla
la scadenza e invia il token solo alle API del progetto. Un 401 su una
richiesta protetta porta al login; un 403 non comporta logout.

I guard e i pulsanti nascosti migliorano l'interfaccia, ma la sicurezza
dipende dalle autorizzazioni backend. ADMIN crea utenti, archivia e marca
duplicati. USER cambia stato soltanto per le proprie assegnazioni;
READONLY non modifica issue. Le notifiche sono consultabili soltanto dal
destinatario, anche rispetto ad ADMIN. Secondo la decisione del gruppo su
RF-17, solo USER e ADMIN possono segnare lette le proprie notifiche; READONLY
puo' soltanto consultarle. Il rifiuto server e i comandi UI coerenti con
questa decisione sono verificati nelle evidenze del 19 ottobre.
Il link Utenti e' mostrato solo ad ADMIN; la rotta resta protetta anche con
navigazione diretta.

Il setup locale usa HTTP e credenziali di sviluppo. Un'eventuale esposizione
pubblica richiederebbe HTTPS e configurazione adeguata dei segreti;
questo documento non certifica un deployment di produzione.
Connessione, chiave JWT e credenziali di bootstrap non sono incorporate nei
file appsettings. Compose passa i valori dimostrativi; l'avvio nativo richiede
variabili d'ambiente o user-secrets. Il seeder non modifica un ADMIN gia'
presente. L'API rifiuta enum numerici non definiti per tipo, priorita' e stato;
la priorita' omessa resta valida.

\section{Ricerca, suggerimento ed export}
La lista esclude gli archiviati e applica filtri per tipo, stato, priorita',
creatore e assegnatario dove previsti dall'API. La ricerca API usa
\texttt{ILike} di PostgreSQL su titolo e descrizione. L'helper
\texttt{SearchIssues} opera invece su collezioni in memoria: i suoi test
non sostituiscono i test HTTP della query SQL.

Il suggerimento considera gli utenti attivi ADMIN e USER. Conta soltanto
issue assegnate non archiviate in TODO o IN\_PROGRESS. Vince il carico
minore; a parita', l'ID minore. Nessun candidato produce 404 con messaggio.
La consultazione del suggerimento non scrive dati.

Lista ed export riutilizzano la stessa query backend. Il CSV contiene
Id, Title, Type, Priority, Status, CreatedByEmail, AssignedToEmail e
CreatedAt; usa UTF-8 ed escaping di virgole, doppi apici e ritorni a capo.
Il frontend scarica il Blob con nome \texttt{issues.csv} e passa i filtri
correnti anche durante l'attesa della ricerca live. RxJS \texttt{switchMap}
annulla le richieste precedenti evitando risultati obsoleti.
Il download browser aggiunge BOM UTF-8 e la dichiarazione \texttt{sep=,}
per l'apertura in Excel; la risposta CSV standard dell'API resta invariata.

\section{Immagini, archivio e duplicati}
L'immagine viene caricata dopo la creazione come multipart con campo
\texttt{image}. Sono ammessi PNG/JPEG non vuoti entro 5 MiB (5242880 byte).
Il backend controlla estensione e marcatori del file, non esegue una
decodifica completa o una scansione antivirus. Solo ADMIN o USER creatore
puo' caricare; tutti i ruoli autenticati possono leggere. L'immagine
viene scaricata tramite HttpClient con JWT e visualizzata con un object URL,
revocato quando la pagina viene distrutta.

Un upload fallito non annulla la creazione: il form conserva l'ID e riprova
solo il caricamento. L'archiviazione imposta \texttt{IsArchived}, senza
cancellazione; l'archivio mantiene accessibili dettaglio e immagine.
Il duplicato imposta DUPLICATE e \texttt{DuplicateOfIssueId}, azzera
\texttt{ResolvedAt} e non genera la notifica di risoluzione.

\section{Diagrammi comportamentali}
Le viste seguenti descrivono creazione, cambio stato e duplicazione.
\diagramma{../Secondo homework/diagrammi_comportamentali/sequence_creazione_issue.png}{Sequenza di creazione della issue.}
\diagramma{../Secondo homework/diagrammi_comportamentali/sequence_cambio_stato.png}{Sequenza di cambio stato e notifica.}
\diagramma{../Secondo homework/diagrammi_comportamentali/sequence_issue_duplicata.png}{Sequenza di contrassegno duplicato.}
\clearpage

\section{Interfacce e avvio}
\begin{longtable}{P{0.12\textwidth} P{0.49\textwidth} P{0.30\textwidth}}
\toprule Metodo & Percorso API & Scopo\\\midrule\endhead
POST & \path{/api/auth/login} & Login.\\
GET / POST & \path{/api/users} & Lista/creazione utenti, solo ADMIN.\\
GET / POST & \path{/api/issues} & Lista/creazione issue.\\
GET & \path{/api/issues/{id}} & Dettaglio.\\
PATCH & \path{/api/issues/{id}/status} & Cambio stato.\\
PATCH & \path{/api/issues/{id}/archive} & Archiviazione ADMIN.\\
GET & \path{/api/issues/archived} & Archivio.\\
PATCH & \path{/api/issues/{id}/duplicate} & Contrassegno duplicato ADMIN.\\
GET & \path{/api/issues/suggest-assignee} & Suggerimento.\\
GET & \path{/api/issues/export} & CSV.\\
POST / GET & \path{/api/issues/{id}/image} & Upload/lettura immagine.\\
GET & \path{/api/notifications} & Notifiche personali.\\
PATCH & \path{/api/notifications/{id}/read} & USER/ADMIN segnano lette le proprie notifiche; READONLY riceve 403.\\
\bottomrule
\end{longtable}
Le rotte frontend sono \path{/login}, \path{/issues}, \path{/issues/new},
\path{/issues/:id}, \path{/issues/archived}, \path{/users} e
\path{/notifications}. La pagina utenti e' riservata ad ADMIN;
la creazione e' protetta anche contro l'accesso diretto READONLY.

Prerequisiti: Docker Desktop con motore Linux, Node.js 22/npm; .NET SDK 10
per i test o per l'avvio backend fuori da Docker. Dalla radice:
\begin{verbatim}
docker compose -f docker/docker-compose.yml up -d --build
# In un altro terminale:
cd frontend
npm ci
npm start
\end{verbatim}
Frontend: \url{http://localhost:4200/login}; API:
\url{http://localhost:5256/api}. Attendere l'avvio delle API prima del login.
Su database nuovo l'ADMIN e' \texttt{admin@bugboard26.local}, password
\texttt{Admin123!}: sono credenziali locali di sviluppo. Il README contiene
le alternative di configurazione. Non usare \texttt{down -v} per conservare
i volumi. L'avvio con \texttt{dotnet run} usa una directory immagini diversa
dal container e non trasferisce automaticamente i file. L'avvio nativo
richiede connessione PostgreSQL, chiave JWT e credenziali di bootstrap
configurate tramite ambiente o User Secrets, come descritto nel README;
non sono conservati segreti di avvio nei file applicativi versionati.

\section{Versionamento e artefatti}
Il repository Git indicato all'inizio raccoglie codice, test e documenti.
La cronologia Git documenta l'evoluzione dei sorgenti e i contributi;
il riferimento della revisione descritta e' riportato in copertina.
Sorgenti PlantUML e PNG correnti sono verificati insieme; i file VPD sono
conservati come artefatti storici. Provenienza delle viste e limiti della
verifica CASE sono descritti nel capitolo 4.

\chapter{Testing e valutazione dell'usabilita'}
\section{Strategia e livelli di verifica}
Le verifiche sono separate in test di dominio, integrazione backend e
collaudo frontend con sistema reale. I test di dominio usano xUnit e
collezioni in memoria. I test HTTP backend usano WebApplicationFactory
con autenticazione reale e PostgreSQL dedicato; non richiedono una porta
API esterna. Ogni caso usa un database temporaneo distinto. Il database
di sviluppo non viene svuotato dai test.

Il frontend usa Jasmine/Karma. I suoi test simulano servizi o HTTP per
isolare il comportamento dei componenti: un test passato non dimostra
da solo il collegamento al backend. Il collaudo del 16 ottobre ha invece
usato Angular, backend Docker e PostgreSQL reali, con Chrome headless.

\section{Tre metodi di dominio e classi di equivalenza}
\begin{longtable}{P{0.29\textwidth} P{0.64\textwidth}}
\toprule Metodo & Casi presenti e motivazione\\\midrule\endhead
\texttt{SearchIssues}\newline(issues, keyword) & Parola nel titolo, nella descrizione e parola assente: distingue i due campi di ricerca e la classe senza corrispondenze.\\
\texttt{CanChangeIssueStatus}\newline(user, issue) & ADMIN non assegnatario; USER assegnatario/non assegnatario; READONLY anche assegnato; USER inattivo. Separa ruolo, assegnazione e attivazione.\\
\texttt{SuggestAssignee}\newline(users, issues) & Carichi differenti con vincitore alternato, issue risolte/chiuse escluse e nessun candidato assegnabile. Verifica la scelta e il caso limite senza risultato.\\
\bottomrule
\end{longtable}
I tre metodi hanno ciascuno due parametri. Le Theory rappresentano casi
significativi rispetto alle regole del dominio. La suite di integrazione aggiunge
verifiche di permessi, filtri, escaping CSV, parita' ed esclusioni del
suggerimento, immagini e persistenza. La copertura strutturale completa
non e' deducibile da questo elenco: va misurata sulla suite completa.
\subsection{Classi coperte e limiti della strategia}
Riferimenti: \path{backend/BugBoard26.Api/Services/IssueDomainService.cs} e
\path{backend/BugBoard26.Tests/IssueDomainServiceTests.cs}. La strategia usa
classi di equivalenza e regole decisionali; non costituisce una suite completa
SECT/R-WECT e non dimostra copertura strutturale completa.

\textbf{SearchIssues.} Le keyword presenti sono \texttt{login}, \texttt{API}
e \texttt{grafica}: coprono titolo, descrizione e assenza di risultati.
Il test asserisce il conteggio, non gli ID restituiti. Keyword nulla, vuota,
di soli spazi, trimming, differenze di maiuscole, elenco vuoto e molteplici
corrispondenze non hanno casi dedicati in questo file. In particolare, la
keyword API e' identica al testo della fixture: non prova da sola il confronto
case-insensitive del metodo. Le prove PostgreSQL della ricerca API restano
distinte da quelle dell'helper in memoria.

\textbf{CanChangeIssueStatus.} Le classi coperte separano ADMIN attivo,
USER assegnatario/non assegnatario, READONLY assegnatario e USER inattivo.
L'oracolo e' il booleano atteso. Non sono coperte tutte le combinazioni:
per esempio ADMIN inattivo e USER attivo con assegnatario nullo non hanno
casi unitari dedicati. Date, transizioni effettive e salvataggi appartengono
al controller e richiedono test HTTP/persistenza, non questo metodo puro.

\textbf{SuggestAssignee.} Le Theory usano carichi (0,2) e (3,1) e verificano
l'ID del vincitore; le fixture escludono un READONLY e issue Resolved/Closed.
Il Fact usa soltanto READONLY attivo e USER inattivo e attende \texttt{null}.
Parita' per ID, ADMIN candidato, InProgress, archiviati ed elenco vuoto non
hanno casi unitari dedicati nel file. Il codice implementa queste regole e
sono presenti scenari HTTP aggiuntivi; non si trasforma la loro presenza
in copertura unitaria misurata. La suite del 19 ottobre esegue anche questi
test HTTP: 166 casi complessivi superati, 11 di dominio e 155 di integrazione.
Il report e il TRX sono in \path{verifica-backend/}; non si attribuiscono
alle sole prove unitarie i risultati ottenuti attraverso l'API.


\section{Test plan -- Risoluzione di una issue}
\textbf{Obiettivo:} verificare UC-02 e RF-10/RF-11 senza consentire
modifiche a utenti non autorizzati. Il piano prevede una issue di prova
creata da S e assegnata ad A, entrambi USER attivi e distinti; U e' un altro
USER, R e' READONLY, con un ADMIN separato. Stato, date e notifiche vengono
confrontati prima e dopo le operazioni, usando dati predisposti per la prova.
\subsection{Ambiente, criteri e raccolta delle evidenze}
Ambito: UC-02, RF-10, RF-11, RF-17, RNF-02 e feedback UI di RNF-05.
Sono esclusi da questo piano export, immagini e archivio, coperti da
verifiche di integrazione separate. L'ambiente comprende .NET 10, xUnit
e PostgreSQL 16 dedicato ai test sulla porta 55432.
Il controllo dal browser richiede inoltre Angular e l'API del setup di prova.
Gli account di prova e gli ID sono predisposti dalle fixture. I casi
di rifiuto confrontano snapshot persistenti prima e dopo la richiesta;
account inattivo e issue duplicata sono condizioni separate.

\textbf{Criteri di ingresso:} disponibilita' di API, login e database isolato,
con data, riferimento dei sorgenti, versioni, ambiente e comando registrati.
L'assenza di un prerequisito comporta un esito bloccato.
La suite HTTP usa database distinti per caso e JWT ottenuti dal login reale.
Una fixture puo' predisporre condizioni non ottenibili dalla UI, per esempio
un READONLY assegnatario, senza modificare i dati di sviluppo.

\textbf{Oracoli:} stato, date UTC, destinatario e notifiche persistenti.
Due chiamate separate a \texttt{UtcNow} non devono necessariamente produrre
date identiche. Stesso stato valido: date e notifiche restano invariate;
rifiuto: nessuna scrittura. La verifica associa alla risposta HTTP il
riscontro persistente e l'evidenza dell'esito. Il controllo del proprietario
si applica anche quando chi tenta la lettura e' ADMIN.

\textbf{Evidenze acquisite:} suite backend, coverage, persistenza e analisi
qualita' sono in \path{verifica-backend/} sotto Documentazione finale.
Il report frontend storico e' in \path{verifica-frontend/}; l'audit del
19 ottobre registra i nuovi esiti. RF-17 e' verificato in backend e UI.
Il CSV associa i casi ai metodi passati nel TRX e ai riscontri browser,
distinti dalle tre sessioni con utenti descritte nella valutazione di usabilita'.

La tabella identifica quattordici casi, TP-01--TP-14. Stato e permessi sono in
\path{IssueStatusIntegrationTests.cs}; duplicati in
\path{IssueDuplicateValidationIntegrationTests.cs}; notifiche in
\path{NotificationsIntegrationTests.cs}. Riapertura/seconda risoluzione,
account disattivato e ripetizione del PATCH read hanno esiti dedicati
nel TRX. Per TP-14, RF-17 fissa l'oracolo: GET delle
proprie notifiche 200; PATCH read su notifica propria esistente 403, anche
se gia' letta, senza modificare \texttt{IsRead}, date o altre righe.
Le fixture includono per R notifiche lette e non lette. Il riscontro browser
verifica l'assenza del comando UI e di PATCH automatici durante la consultazione;
USER e ADMIN mantengono la lettura delle proprie notifiche (TP-13 e regressione
per ADMIN). Esiti e riferimenti sono nel CSV; il criterio RF-17 e' invariato.


\begin{longtable}{P{0.09\textwidth} P{0.38\textwidth} P{0.43\textwidth}}
\toprule Caso & Azione & Risultato atteso\\\midrule\endhead
TP-01 & A imposta RESOLVED. & 200; stato e date persistiti; nuova notifica non letta per S.\\
TP-02 & ADMIN non assegnatario aggiorna una issue di prova. & Operazione consentita.\\
TP-03 & U tenta il cambio stato. & Nessun pannello frontend; PATCH diretto 403 senza scritture.\\
TP-04 & R tenta il cambio stato. & Nessun pannello; PATCH diretto 403.\\
TP-05 & Richiesta senza JWT. & 401 e nessuna modifica.\\
TP-06 & ID issue inesistente. & 404 e nessuna notifica.\\
TP-07 & A ripete lo stato RESOLVED corrente. & Nessuna nuova notifica o variazione delle date.\\
TP-08 & Riapertura e seconda risoluzione. & Data azzerata alla riapertura; nuova data/notifica alla seconda risoluzione.\\
TP-09 & Cambio stato su duplicata. & 409; riferimento e stato conservati.\\
TP-10 & S consulta notifiche; U tenta lettura della notifica di S. & S vede la propria notifica; U non la vede e PATCH read altrui restituisce 403.\\
TP-11 & JWT di account diventato inattivo. & 401 dal lookup dell'account; nessuna modifica persistente.\\
TP-12 & Si richiede Duplicate nel percorso status. & 400; stato, riferimento, date e notifiche invariati.\\
TP-13 & S segna letta la notifica e ripete il PATCH. & 200; lettura persistita; nessuna notifica aggiunta; ID assente 404, senza token 401.\\
TP-14 & READONLY consulta le proprie notifiche e tenta di segnarne una letta. & GET 200, solo notifiche proprie; PATCH read 403 e dati invariati. UI senza comando di lettura. Decisione del gruppo RF-17 verificata.\\
\bottomrule
\end{longtable}
\textbf{Criterio di accettazione:} tutti i casi pianificati eseguiti con
esito registrato e nessuna modifica non autorizzata. Gli esiti effettivi sono
in \path{esiti_test_plan.csv}, con rimandi a metodi di test e TRX; i controlli
UI richiamano il collaudo browser. TP-08 e TP-11 hanno casi HTTP dedicati;
TP-13 include il PATCH ripetuto e TP-14 il rifiuto READONLY anche su notifiche
gia' lette. Gli esiti dei test tecnici sono distinti dalle sessioni di
usabilita' con persone.

\section{Collaudo frontend reale -- 16 ottobre 2026}
Il report del collaudo e gli screenshot sono in
\path{verifica-frontend/report.md} e \path{verifica-frontend/immagini/}.
Riferimento: commit \texttt{e13995140dff382dc3d548208c6c1e60be91fe5d},
checkout pulito prima e dopo. API sulla porta 5256 e Angular sulla 4200;
PostgreSQL del Compose di sviluppo. Le prove hanno utilizzato il sistema
avviato, senza sostituire le risposte delle API con simulazioni.

Sono stati creati tre account con prefisso
\texttt{collaudo-20261016084339} (due USER e un READONLY) e le issue 4--7.
Le operazioni di modifica hanno riguardato soltanto queste nuove issue.
La issue 4 e' stata duplicata e archiviata; la 7, predisposta tramite API
per avere un assegnatario deterministico, e' stata risolta dal relativo USER.
Le altre creazioni sono passate dal form frontend. Non sono riportati
token o credenziali dei partecipanti nei risultati.

\begin{longtable}{P{0.32\textwidth} P{0.61\textwidth}}
\toprule Area & Esito osservato\\\midrule\endhead
Login e guard & Login ADMIN riuscito; password errata: 401 e messaggio UI. Accesso anonimo a creazione reindirizzato al login; nessuna navbar nella login.\\
Ricerca e lista & Ricerca live, digitazione rapida con risultato finale corretto, filtri tipo/stato/priorita', ordinamento per titolo, lista vuota e reset completo riusciti.\\
Export & Download reale; ID e ordine nel CSV uguali alla lista filtrata.\\
Creazione & Issue senza immagine e con PNG create; immagine scaricata autenticata e visualizzata.\\
Recupero upload & JPEG volutamente non valido rifiutato dall'API; ID conservato; JPEG valido caricato al secondo tentativo senza seconda creazione.\\
Validazione & Titolo/descrizione di soli spazi: nessuna POST e messaggi sui campi. Formato non ammesso e file oltre 5 MiB rifiutati dal frontend.\\
Stato e ruoli & USER assegnatario risolve; stato persistito. USER non assegnatario senza pannello. USER non assegnatario e READONLY ricevono 403 dal PATCH diretto.\\
READONLY & Nessuna azione di modifica nel dettaglio; accesso diretto a creazione reindirizzato; immagine consultabile.\\
Duplicato e archivio & ADMIN marca duplicato con collegamento aggiornato; archivia; issue assente dalla lista principale e dettaglio apribile dall'archivio.\\
Errori e layout & ID inesistente mostra ``Issue non trovata''. Lista, form e dettaglio controllati a 1366 e 390 pixel senza overflow orizzontale; screenshot ispezionati.\\
Notifica & Via API verificata la notifica di risoluzione della issue 7 al segnalatore ADMIN; non un collaudo completo della pagina notifiche.\\
\bottomrule
\end{longtable}
La build frontend e' riuscita e i 27 test Jasmine/Karma sono passati senza
watch con ChromeHeadless. La prova non ha richiesto correzioni applicative.
Gli screenshot selezionati sono conservati in
\path{verifica-frontend/immagini/}.

Questa prova non misura le prestazioni con rete rallentata o grandi
quantita' di dati, non copre ogni browser e non sostituisce il collaudo
di persistenza dopo ricreazione dei container.

\section{Esecuzione ripetibile e qualita' backend}
Le verifiche successive sono registrate separatamente in
\path{verifica-backend/} e \path{verifica-frontend/review-20261019/}.
I report includono comandi, data, HEAD, modifiche non committate e SHA256.
Il 19 ottobre sono stati eseguiti 166 casi backend (11 dominio, 155 integrazione)
con PostgreSQL dedicato: tutti superati, nessuno ignorato. Il test plan e'
ricondotto ai singoli metodi nel CSV degli esiti.

Il collaudo integrato dell'audit del 19 ottobre ha verificato 20 controlli con API Docker
reali e database E2E dedicato, senza mock HTTP: utenti, creazione con PNG,
suggerimento, dettaglio, risoluzione, notifiche USER/ADMIN, RF-17, CSV,
duplicazione e archivio. Le correzioni UI successive riguardano la visibilita'
del link Utenti e il nome accessibile della ricerca; il report di audit
registra la verifica successiva. Le prove con utenti restano distinte.

Dalla radice del repository, con Docker disponibile:
\par
\begin{minipage}{\textwidth}
\begin{verbatim}
docker compose -p bugboard26-tests `
  -f backend/BugBoard26.Tests/docker-compose.tests.yml up -d --wait
dotnet test backend/BugBoard26.slnx
# Frontend, dalla cartella frontend:
npm run build
node node_modules/@angular/cli/bin/ng.js test `
  --watch=false --browsers=ChromeHeadless
\end{verbatim}
\end{minipage}\par
Il database di test usa la porta 55432 ed e' distinto da quello di sviluppo.
La sola esecuzione dei test di dominio con filtro non sostituisce la suite
di integrazione. Dopo l'adattamento del download CSV i test Jasmine/Karma
sono 40, inclusi controlli di navigazione e creazione per ruolo e verifica
dei byte esportati, eseguiti in ChromeHeadless senza watch. La build e'
riuscita e il browser reale ha verificato il download. Il resoconto del
collaudo interno riporta inoltre l'apertura corretta del CSV in Excel.

SonarQube locale e la raccolta OpenCover sono predisposti in
\path{quality/analyze-backend.ps1} e \path{quality/coverage.runsettings}.
L'analisi corrente e' stata elaborata dal server locale Community Build
26.10.0.133078, scanner .NET 11.3.0: coverage importata per 29 file principali,
Quality Gate OK con condizioni valutate, zero bug e zero vulnerabilita'
segnalati, quattro code smell lasciati visibili. Coverage Sonar complessiva
93,7\%, linee 95,7\%, rami 88,2\%. Non si dichiara copertura totale o assenza
di ogni possibile difetto. Gate e percentuali sono riferiti a HEAD con le
correzioni locali, non al solo commit \texttt{6aa175a}.

Tre code smell S6932 riguardano la lettura manuale del multipart, mantenuta
per preservare autorizzazioni e risposte 400 uniformi; un avviso ASP0027
riguarda Program pubblico parziale. Non sono introdotte esclusioni o
soppressioni. L'avviso storico su Microsoft.OpenApi 2.0.0 e' superato dal pin
2.7.5. Dati e immagini sono stati confrontati dopo restart e ricreazione del
solo backend, senza eliminare volumi; il rebuild successivo conserva la
stessa issue e lo stesso SHA256 immagine. Nessun test automatico usa il
database di sviluppo; le prove Docker con dati riconoscibili sono distinte.

\section{Valutazione dell'usabilita' -- review e sessioni}
La valutazione comprende una review tecnica sistematica e tre sessioni
con partecipanti potenziali, documentate dal resoconto del moderatore e
dal questionario. I due metodi sono complementari: il collaudo tecnico
verifica comportamenti del sistema, mentre le sessioni raccolgono
riscontri d'uso. I soli voti non permettono di ricostruire tempi,
attivita' non osservate o cause delle difficolta'.

\subsection{Review tramite checklist}
La review sistematica assistita del 19 ottobre e' registrata in
\path{evidenze/review_tecnica.csv}, con otto criteri, evidenze e limiti.
Il report completo e' in \path{verifica-frontend/review-20261019/report.md};
la prima review backend resta storica. Sono ispezionati sorgenti e screenshot
e provati 28 controlli browser con API Docker reale e PostgreSQL dedicato.
La review e' assistita da automazione browser e non costituisce una
valutazione umana indipendente o una certificazione WCAG.
\begin{itemize}
\item Terminologia coerente fra lista, form e dettaglio.
\item Navigazione riconoscibile e ritorno alla lista disponibile.
\item Campi obbligatori ed errori comprensibili vicino al campo.
\item Caricamento, esito del salvataggio, errore e lista vuota distinguibili.
\item Azioni coerenti con il ruolo, senza pulsanti di modifica per READONLY.
\item Riconoscibilita' della notifica di risoluzione e dello stato letto/non letto.
\item Uso da tastiera, focus visibile e leggibilita' su schermi piccoli.
\item Recupero dell'upload senza indurre una seconda creazione.
\end{itemize}

Login, ricerca e creazione sono percorsi con Tab, Shift+Tab ed Enter,
controllando il focus visibile. Sono distinti caricamento, lista vuota e
errore di rete; ritardo e interruzione sono simulati soltanto nel browser,
con recupero tramite Pulisci. Un falso PNG e' invece realmente rifiutato
dall'API: il form mostra l'ID gia' salvato e il nuovo upload valido usa lo
stesso ID, con una sola richiesta di creazione. Nessun dato di sviluppo
e' modificato. Quattro viste sono controllate a 320 e 390 pixel senza
overflow orizzontale rilevato.

Quattro dei cinque campioni di contrasto della lista erano insufficienti
e sono corretti: intestazioni e nota export passano a 4,97:1, priorita'
a 5,25:1, stato a 5,63:1; il pulsante resta a 4,70:1. Le soglie sono
confrontate sui valori non arrotondati. Sono scurite anche note analoghe
del form/dettaglio e testo di errore login, senza ridisegnare l'interfaccia.
I 40 test frontend esistenti e la build sono superati nuovamente; non sono
aggiunti test permanenti. Restano migliorabili annunci per lettori schermo,
focus iniziale sugli errori e conferma delle azioni admin. La review e'
completata nel perimetro dichiarato, non e' una certificazione di accessibilita'.

\subsection{Sessioni con utenti e questionario}
Il protocollo definisce cinque attivita': ricerca di una issue (U1),
creazione (U2), risoluzione come assegnatario (U3), consultazione e lettura
della notifica (U4), filtro ed export in sola lettura (U5).
Nelle tre sessioni registrate sono state svolte soltanto U2 e U4;
non si attribuisce ai partecipanti l'esecuzione dell'intero protocollo.

Il questionario usa una scala da 1 (per nulla) a 5 (molto) per navigazione,
comprensibilita' dei campi, ricerca, esiti delle operazioni e permessi.
Prevede inoltre una domanda sul passaggio meno chiaro e un suggerimento
facoltativo. Per queste due domande non sono disponibili risposte.
Si tratta di un questionario locale breve, non di uno strumento SUS.
Protocollo e testi delle attivita' sono in \path{usabilita.md}; sessioni e
punteggi sono conservati in \path{evidenze/sessioni_utenti.csv} e
\path{evidenze/questionario.csv}.

\subsection{Risultati e limiti delle tre sessioni}
Fonte: resoconto del moderatore e risposte trascritte del questionario;
non sono allegate registrazioni delle sessioni. Data: 19 ottobre 2026.
Il campione di convenienza comprende tre persone conosciute dal moderatore,
che hanno usato il suo PC in condivisione, tutte con ruolo USER. P01 e P03
hanno svolto U2, creazione issue; P02 U4, riscontro e lettura notifica.
Per tutti il completamento e' dichiarato autonomo dopo una spiegazione
iniziale, che viene esplicitamente registrata. Non e' descritto il contenuto
dell'introduzione o il mezzo di condivisione.

\begin{center}
\begin{tabular}{llrrrrr}
\toprule Codice & Attivita' & Q1 & Q2 & Q3 & Q4 & Q5\\\midrule
P01 & U2 & 5 & 4 & 4 & 4 & 5\\
P02 & U4 & 4 & 4 & 5 & 4 & 4\\
P03 & U2 & 4 & 3 & 4 & 4 & 3\\
\bottomrule
\end{tabular}
\end{center}
Q1 indica navigazione, Q2 campi, Q3 ricerca, Q4 esiti, Q5 permessi;
scala da 1 (per nulla) a 5 (molto). Sono riportati i valori ricevuti senza
trasformarli in un punteggio SUS o in una validazione statistica.
I voti 3 di P03 per campi e permessi indicano una valutazione moderata:
in assenza di osservazioni non si attribuiscono cause specifiche.

I partecipanti sono identificati con codici senza nomi nei file consegnati;
il moderatore li conosce, quindi non si dichiara anonimato assoluto.
Campione di convenienza ridotto, attivita' diverse e solo ruolo USER limitano
la generalizzabilita'. Non si assume che ciascuno abbia svolto ricerca,
filtri o prove con gli altri ruoli. Profilo, esperienza, browser, checkout
effettivo, tempi e risposte aperte Q6/Q7 non sono comunicati. I campi vuoti
significano non fornito, non assenza di problemi. La modalita' del consenso
non e' documentata. La review tecnica resta un'attivita' distinta.
Dettagli e provenienza: \path{evidenze/risposte-parziali-usabilita.md}.

\chapter{Limiti e conclusioni}
\section{Limiti della valutazione}
Le verifiche tecniche sono riferite alle revisioni e agli ambienti
identificati nei report. Non coprono tutti i browser, condizioni di carico
elevato o un deployment pubblico. Il superamento dei test e il Quality Gate
non rappresentano una certificazione completa di sicurezza o prestazioni.
Le sessioni di usabilita' sono limitate a tre partecipanti, al ruolo USER
e a due attivita'; le lacune del resoconto sono indicate nel capitolo 3.

\section{Provenienza e verifica dei diagrammi}
I materiali storici dei due homework conservati nel repository comprendono
otto archivi VPD con metadati di Visual Paradigm Online. Le viste riportate nella relazione sono
mantenute in sorgenti PlantUML modificabili e corrispondenti immagini PNG.
I modelli descrivono casi d'uso, componenti, deployment locale, classi,
schema dati e tre sequenze applicative.

Il confronto con il codice riguarda entita', relazioni e vincoli di
persistenza, dipendenze dei controller, permessi e flussi delle API.
L'aggiornamento del 19 ottobre allinea in particolare le notifiche personali
e il divieto di scrittura READONLY. Il deployment locale rappresenta la
configurazione Compose con i due volumi persistenti.

Le dodici viste PNG sono verificate rispetto ai sorgenti incorporati nei
metadati delle immagini. Il file \path{diagrammi_manifest.json} conserva
percorsi e impronte SHA256 di sorgenti ed export. La generazione aggiornata
dei casi d'uso e del cambio stato utilizza PlantUML 1.2026.8 e Java 25.
Il controllo riguarda le viste PlantUML correnti e l'integrita' ZIP dei
VPD storici, non l'apertura o la sincronizzazione di questi ultimi
nell'editor originale.

\section{Perimetro del prodotto}
Non sono previsti paginazione, ripristino archivio, refresh token o gestione
immagini dal dettaglio frontend.
Queste esclusioni mantengono il prodotto nel perimetro assegnato.

\section{Conclusioni}
BugBoard26 realizza le funzionalita' assegnate con frontend e backend
indipendenti, dati centralizzati e operazioni regolate dai ruoli.
Test di dominio, integrazione, collaudi browser e analisi statica
supportano i comportamenti verificati. La review ha portato a correzioni
puntuali dell'interfaccia e le sessioni con partecipanti hanno fornito un
primo riscontro d'uso. Questi risultati restano circoscritti ai metodi,
agli scenari e ai limiti dichiarati nella relazione.

\end{document}
